\section{Discuss\~ao}
\label{sec:discussao}

\subsection{O que os dois m\'etodos oferecem, pelo que foi medido}

\begin{center}
\begin{tabular}{p{4.6cm}p{4.6cm}p{4.6cm}}
\toprule
 & front-tracking (balanceado) & VOF \\
\midrule
curvatura / salto de press\~ao & exata no c\'irculo; erro \num{1.7e-7}
  & derivada de campo suavizado; erro \num{4.7e-3} \\
correntes par\'asitas & $<\num{5e-11}$, decaindo & \num{5.2e-5}, decaindo \\
massa & conservada enquanto o escoamento \'e est\'atico; deriva com o movimento
  & conservada por constru\c{c}\~ao, sempre \\
topologia (coalesc\^encia, quebra) & manual, fora do escopo & autom\'atica \\
custo de estrutura & marcadores ordenados, cirurgia, indicadora por c\'elula
  & campo euleriano a mais \\
\bottomrule
\end{tabular}
\end{center}

A leitura n\~ao \'e ``um vence''. Nas m\'etricas deste teste o front-tracking
balanceado sai \`a frente, mas o teste \'e uma gota est\'atica: n\~ao exercita
topologia, n\~ao exercita deforma\c{c}\~ao grande, e n\~ao exercita a
conserva\c{c}\~ao de massa sob movimento --- que \'e justamente onde o VOF tem
vantagem \emph{estrutural}, e n\~ao circunstancial. Escoamento em que a massa da
fase precisa ser conservada ao longo de muitos tempos, ou em que gotas coalescem
e se partem, continua favorecendo o VOF. Com os dois no mesmo sistema, a escolha
passa a ser informada por medida em vez de por prefer\^encia.

\subsection{A compara\c{c}\~ao como verifica\c{c}\~ao m\'utua}

Vale sublinhar um uso da compara\c{c}\~ao que n\~ao \'e o de escolher m\'etodo.
Como a dire\c{c}\~ao de v\'arios resultados \'e conhecida \emph{a priori} pela
literatura, eles funcionam como or\'aculo um para o outro: se o VOF tivesse
sa\'ido melhor em curvatura, ou se o front-tracking espalhado tivesse conservado
massa exatamente, o resultado apontaria defeito de implementa\c{c}\~ao, n\~ao
descoberta sobre os m\'etodos.

E foi exatamente assim que o defeito apareceu. O \'unico resultado que
\emph{contrariou} a expectativa --- correntes par\'asitas maiores no
front-tracking --- n\~ao era propriedade do m\'etodo: era uma deficiência da
implementa\c{c}\~ao, localizada pela discrep\^ancia e removida na
Se\c{c}\~ao~\ref{sec:balanco}. Sem o VOF ao lado, com o mesmo problema e as
mesmas sondas, n\~ao haveria com o que estranhar \num{5e-4}: pareceria apenas
``pequeno''.

\subsection{Limita\c{c}\~oes}

Para que o alcance dos n\'umeros n\~ao seja superestimado:

\begin{itemize}
  \item \textbf{Uma \'unica resolu\c{c}\~ao.} Tudo foi medido a $10{,}2$
    c\'elulas por raio. Um estudo de converg\^encia --- refinar $h$ e obter a
    taxa de cada m\'etrica --- n\~ao foi feito.
  \item \textbf{A vantagem do balan\c{c}o \'e espec\'ifica do caso est\'atico
    com $\kappa$ constante}, como a Se\c{c}\~ao~\ref{sec:elipse} mostrou: na
    elipse as duas formula\c{c}\~oes ficam equivalentes dentro da medida.
  \item \textbf{A curvatura na faceta \'e uma m\'edia ponderada dos marcadores
    vizinhos}, e n\~ao um campo de curvatura propriamente constru\'ido e
    estendido. Foi o suficiente para estabilizar a elipse, mas n\~ao foi testada
    em geometria com curvatura fortemente vari\'avel ou com mudan\c{c}a de
    sinal.
  \item \textbf{Serial --- e agora medido, n\~ao suposto.} Tudo neste
    relat\'orio \'e $n_p=1$. O que mudou \'e que $n_p=2$ foi \emph{testado} nos
    dois m\'etodos, e a subse\c{c}\~ao~\ref{sec:paralelo} traz o resultado: nenhum
    dos dois produz a resposta certa, e nenhum dos dois falha.
  \item \textbf{Propriedades iguais.} A terceira mudan\c{c}a estrutural --- o
    salto de $\rho$ e $\mu$ --- est\'a adiada nos dois lados. Em escoamento com
    raz\~ao de densidade alta as conclus\~oes podem mudar.
  \item \textbf{A cirurgia foi exercitada} apenas na relaxa\c{c}\~ao da elipse,
    onde a frente encurta monotonicamente (108 para 102 marcadores). Esticamento
    forte, que exigiria \emph{inser\c{c}\~ao} sustentada, segue sem teste
    acoplado.
\end{itemize}

\subsection{Paralelo: o que foi medido em $n_p=2$}
\label{sec:paralelo}

Este relat\'orio dizia, at\'e aqui, apenas que o trabalho era serial. Isso era
verdade e insuficiente: ``n\~ao implementado'' e ``implementado e errado'' pedem
provid\^encias diferentes, e a diferen\c{c}a s\'o se conhece rodando. Rodou-se,
em \'arvore de trabalho isolada, o mesmo caso curto com $n_p=1$ e $n_p=2$ nos
dois m\'etodos.

\pgfplotstabletypeset[
  string type,
  columns/grandeza/.style={column name={grandeza}},
  columns/np1/.style={column name={$n_p=1$}},
  columns/np2/.style={column name={$n_p=2$}},
  every head row/.style={before row=\toprule, after row=\midrule},
  every last row/.style={after row=\bottomrule},
]{\dat{paralelo.dat}}

\paragraph{O front-tracking congela quase metade da frente.} Em $n_p=2$,
\textbf{109 dos 252 marcadores n\~ao se moveram} em 10 passos, e formam uma faixa
\emph{cont\'igua} em $y$ --- a parte da bolha que cai no subdom\'inio do outro
processo. A causa \'e uma linha da interpola\c{c}\~ao de velocidade, que em s\'erie
nunca dispara:
\begin{center}
\texttt{if (h <= 0.0) return;\qquad // fora do dominio: marcador nao anda}
\end{center}
A frente \'e \emph{replicada}: cada processo advecta a sua pr\'opria c\'opia com
a velocidade que enxerga, e as c\'opias divergem no primeiro passo. A prova
direta \'e que os dois processos relatam desvios de parti\c{c}\~ao da unidade
\emph{diferentes} ($0{,}402$ e $0{,}795$) --- n\~ao est\~ao vendo o mesmo
est\^encil. Contra \num{3.05e-14} em s\'erie, \'e a mesma assinatura que
identificou a ``for\c{c}a exatamente pela metade'' no corpo r\'igido: pesos que
n\~ao somam 1 devolvem velocidade pequena demais, sem erro vis\'ivel.

\paragraph{O VOF acerta o solver e erra o remalhamento.} Sem remalhar, as
s\'eries de $n_p=1$ e $n_p=2$ saem \textbf{id\^enticas em todos os d\'igitos
impressos}: o caminho multif\'asico do solver e as m\'etricas (que reduzem com
\texttt{MPI\_Allreduce}) est\~ao corretos em paralelo. Com remalhamento, n\~ao:
o pr\'oprio contador do remalhamento acusa \textbf{265 posi\c{c}\~oes sem valor}
no primeiro remalhamento e \textbf{2709} no segundo, e $V_c$ sai com
\SI{77}{\percent} de erro --- \emph{caindo} quando devia subir.

Vale notar o que esse contador \'e: o n\'umero de posi\c{c}\~oes da malha nova
para as quais a colheita n\~ao tinha valor. A mensagem que o imprime existia nos
dois exemplos e vinha imprimindo zero desde sempre, porque em s\'erie \'e zero.
Era um or\'aculo instalado que nunca havia disparado.

\paragraph{Os quatro consertos, e a ordem deles.} Dois s\~ao pequenos e s\~ao o
que falta ao VOF: os escritores de crit\'erio fazem \texttt{fopen(...,"w")} do
\emph{mesmo} caminho em todos os processos, sem guarda de \emph{rank} e sem
barreira antes da releitura. O terceiro n\~ao se cura com guarda: as sementes do
VOF saem do dom\'inio \emph{local} sem \emph{gather}, de modo que mesmo o
processo 0 escreveria uma malha refinada s\'o em torno do peda\c{c}o de interface
que ele possui --- e a fun\c{c}\~ao que resolveria isso j\'a existe em
\texttt{malha-adaptativa.c}, por\'em \texttt{static} e sem uso. (No
front-tracking esse terceiro n\~ao aparece: a frente \'e replicada, ent\~ao todo
processo tem a geometria completa. A assimetria \'e o inverso da da
advec\c{c}\~ao.) O quarto \'e o \'unico sem c\'odigo pronto para copiar: a
advec\c{c}\~ao precisa de velocidade \emph{global} nos marcadores. O
espalhamento, por sua vez, precisa da redu\c{c}\~ao franja$\to$dono que
\texttt{fi\_espalha} j\'a faz vinte linhas adiante, no mesmo arquivo.

\paragraph{Uma armadilha de medi\c{c}\~ao, de gra\c{c}a.} O caso com remalhamento
rodou em \SI{64}{\second} em $n_p=2$ contra \SI{170}{\second} em $n_p=1$: $2{,}7$
vezes em dois processos, \emph{superlinear}. A superlinearidade \'e a assinatura
de resolver um problema \emph{menor}, e n\~ao mais r\'apido. Quem medisse
escalabilidade neste estado veria um resultado excelente.

\paragraph{E o que os tr\^es testes t\^em em comum.} Nenhum quebrou. Tr\^es
execu\c{c}\~oes com c\'odigo de sa\'ida zero, resultados de aspecto plaus\'ivel,
e defeitos de \SI{43}{\percent}, \SI{77}{\percent} e \SI{100}{\percent}. O que os
separou do correto foram os or\'aculos --- um deles j\'a instalado e nunca
disparado.

\subsection{Pr\'oximos passos}

Em ordem de valor:

\begin{enumerate}
  \item \textbf{Converg\^encia em $h$} das tr\^es m\'etricas, nos dois m\'etodos.
    D\'a \`a compara\c{c}\~ao o alcance que hoje lhe falta.
  \item \textbf{Fase B3}, oscila\c{c}\~ao de gota contra a frequ\^encia de Lamb:
    exercita din\^amica acoplada e, por deformar a interface, finalmente p\~oe a
    cirurgia \`a prova com o solver no circuito.
  \item \textbf{Paralelo.} A subse\c{c}\~ao~\ref{sec:paralelo} localiza os
    quatro consertos e os ordena: dois deles s\~ao pequenos e habilitam o VOF;
    a advec\c{c}\~ao global do front-tracking \'e o \'unico que n\~ao tem
    c\'odigo pronto para copiar.
  \item \textbf{Fechar o Hysing}: a Se\c{c}\~ao~\ref{sec:hysing} monta a
    compara\c{c}\~ao contra os valores publicados e declara de antem\~ao como
    ler cada desfecho, mas as corridas at\'e $t=3$ n\~ao terminaram a tempo
    deste texto -- o fechamento segue pendente, e \'e o \'unico item desta
    lista que j\'a est\'a em execu\c{c}\~ao.
\end{enumerate}

A compara\c{c}\~ao com o VOF pode acompanhar cada uma dessas etapas: o caso
existe, as sondas est\~ao instaladas nos dois lados, e as duas formula\c{c}\~oes
de for\c{c}a coexistem atr\'as de um seletor.
